\section{Joint construction of the continuum through compatible histories and prolongations}
\label{act21:iv:continuo-conjunto}
The arithmetic expressed by a normal form preserves a precise reading of the state that produces it. To situate that reading within HMT, it is also necessary to retain the histories, the sheets, and the operations acting upon them. This chapter brings together that common substrate. Its different constructions are obtained from the same tree of APP–TRIT–TPK executions; the expository divisions allow their properties to be studied without turning them into independent domains chosen by a subsequent application.

The starting point is the formal kernel presented above. An admissible emission contains the APP operation, the TRIT orientation, the change of sheet, the residue and the quotient, together with the effective TPK memory update. When only a residue is expressed, the execution history still belongs to the source object. Two executions with an identical final reading are not identified for that reason alone.

\subsection{Admissible histories and memory by composition}
Let \(H_0:=\{\ast\}\). The level \(H_1\) is the finite nonempty set of
admissible APP--TRIT seeds, and, for \(k\geq1\), \(H_k\) is the set of
enriched histories of length \(k\) admitting compatible TPK prolongations
at every finite horizon. Its elements are constructed through successive emissions
from the initial domain of the kernel. We write \(\zeta\varepsilon\)
for the prolongation of \(\zeta\) by an admissible emission \(\varepsilon\), and
\(\rho_k(\zeta\varepsilon)=\zeta\) for deletion of the final emission. The label
produced by the emission is retained, even when different labels express
the same scalar state. The space of complete histories is defined by
\[
 H_\infty^{\rm enr}:=\varprojlim_k(H_k,\rho_k).
\]
The cofinal subsystem \(X_N:=H_{6N}\), with \(X_0=H_0\), coincides with the
sextet clock used in the cardinal proof; by cofinality there is a canonical
identification
\(H_\infty^{\rm enr}\cong\varprojlim_N(X_N,\rho_{6(N+1),6N})\).

On the two arithmetic sheets, the local data satisfy the exact divisions

\[
\delta_\diamond=r_\diamond+9q_\diamond,
\qquad 0\leq r_\diamond<9,
\qquad \diamond\in\{\Sigma,\Pi\}.
\]
The oriented TRIT coordinate similarly admits an expression \(\Delta=o_r+3\kappa\), with the oriented representative declared by the kernel. Preservation of the quotient allows these representations to be composed without replacing an emission by its isolated residue.

In a memory chart, the update of an emission is written

\[
U_\varepsilon(m)=C_\varepsilon m+\kappa_\varepsilon.
\]
For a path \(\alpha=(\varepsilon_1,\ldots,\varepsilon_k)\), the effective operation is the ordered composition of those updates. Its linear part and its translation are

\[
C_\alpha=C_{\varepsilon_k}\cdots C_{\varepsilon_1},
\qquad
\kappa_\alpha=
\sum_{j=1}^{k}C_{\varepsilon_k}\cdots
C_{\varepsilon_{j+1}}\kappa_{\varepsilon_j}.
\]
The empty product in the last sum is the identity.

\begin{proposition}[Affine compatibility of concatenation]
\label{act21:prop:continuo-compatibilidad-afin}
If \(\alpha\) precedes \(\beta\), then

\[
C_{\beta\alpha}=C_\beta C_\alpha,
\qquad
\kappa_{\beta\alpha}=C_\beta\kappa_\alpha+\kappa_\beta,
\qquad
U_{\beta\alpha}=U_\beta U_\alpha.
\]
\end{proposition}
\begin{proof}
Substituting \(U_\alpha(m)\) into \(U_\beta\) yields \(C_\beta C_\alpha m+C_\beta\kappa_\alpha+\kappa_\beta\). Decomposition into linear part and translation gives the three identities. Associativity of composition proves that the result does not depend on an auxiliary subdivision of the path.
\end{proof}

This cocycle identity is the precise content of the accumulated memory. The quantity transported by a second path depends on the chart update it performs; in general, it is not the untransported sum of two records.

The nonadic return preserves a second distinction. In the phase and turn coordinates, with \(\bar r\in\{0,\ldots,8\}\), its representation is

\[
\sigma_9(w,r)=
\left(w+\left\lfloor\frac{\bar r+1}{9}\right\rfloor,
r+1\pmod9\right),
\qquad
\sigma_9^9(w,r)=(w+1,r).
\]
Over nine steps, exactly one crossing from \(8\) to \(0\) occurs. The phase returns and the turn coordinate increases. The enriched holonomy also retains the memory updates accompanying that traversal.

\subsection{Tree refinement and measure of histories}
The initial measure is the canonical one from the census: \(\mu(\ast)=1\) and
\(\mu(c)=1/|H_1|\) for each seed \(c\in H_1\). In this subsection, each
history has a finite positive number of prolongations. Deletion of
emissions defines the refinement maps. Admissibility of the kernel
provides the prolongations; when a residence uses a neutral emission,
that emission explicitly provides a section of \(\rho_k\). The existence of this
section concerns the tree that includes that emission, not an arbitrary
restriction of routes.

Let \(d(\zeta)\geq1\) be the number of admissible emissions at \(\zeta\). The
APP order and the census of prolongations determine the canonical law
\(p_{\rm can}(\varepsilon\mid\zeta)=1/d(\zeta)\). Without additional
parameters, this law defines

\[
\mu(\zeta\varepsilon)=\frac{\mu(\zeta)}{d(\zeta)}.
\]
A nonuniform weighting may replace \(p_{\rm can}\) only when it has already
been produced and added to the enriched record. The correlative constructor
of this chapter uses \(p_{\rm can}\) exclusively; therefore, its measure
is a function of the APP--TRIT--TPK tree and not of an external choice.

\begin{proposition}[Cylindrical compatibility and conditional expectation]
\label{act21:prop:continuo-esperanza-condicional}
At every level,

\[
\sum_{\rho_k(y)=x}\mu(y)=\mu(x).
\]
In the spaces \(L^2(H_k,\mu_k)\), the maps

\[
(\mathsf J_kf)(y)=f(\rho_k(y)),
\qquad
(E_kg)(x)=\sum_{\rho_k(y)=x}\frac{\mu(y)}{\mu(x)}g(y)
\]
satisfy \(E_k=\mathsf J_k^*\), \(E_k\mathsf J_k=I\), and \(\|\mathsf J_kf\|=\|f\|\).
\end{proposition}
\begin{proof}
The first identity is the normalization of the local probabilities. Applying it to \(|f(x)|^2\) yields the isometry. Grouping the sum by fibers of \(\rho_k\) yields the formula for \(E_k\); on a function constant on each fiber, that formula returns its value.
\end{proof}

The compatible masses define a probability measure on the space of infinite histories: the cylinder of a finite history receives its mass \(\mu(\zeta)\). The measure extends from the algebra of cylinders, and its uniqueness follows because they generate the \(\sigma\)-algebra. If the tree is finitely branching and every history has a prolongation, the inverse limit of its levels under deletion of emissions is nonempty. In the cylinder topology, every nonempty cylinder has positive measure.

This measure distinguishes histories converging to the same representation. The uniform measure on terminal values would be a different reading, since it would eliminate the multiplicities that the construction has just preserved.

\subsubsection{Enriched readers and fiber memory}
\label{act21:sec:ch-memoria-fibra}

The difference between history and value admits an operational criterion. At a
finite level \(H_k\), let \(q_k:H_k\to Z_k\) be a reader and let
\(M_k:H_k\to\mathcal M_k\) be the preserved record. The enriched reader
is \(\widehat q_k=(q_k,M_k)\). Memory distinguishes states within a single
fiber of the reader; its sufficiency is verified on the domain under consideration.
This formulation incorporates into the article itself the complete development of the
fiber information used below.

\begin{proposition}[Recoverability at a finite level]
\label{act21:prop:ch-recuperabilidad-fibra}
An inverse of \(\widehat q_k\) exists on its image if and only if
\(\widehat q_k\) is injective. For a random variable \(Z\) taking
values in \(H_k\) with distribution \(\mu_k\),
\[
 \mathsf H(Z)=\mathsf H(q_k(Z))+\mathsf H(Z\mid q_k(Z)).
\]
If \(\widehat q_k\) is injective on the support of \(Z\), then
\[
 \mathsf H(Z\mid q_k(Z),M_k(Z))=0.
\]
Here \(\mathsf H(Z)=-\sum_z\Pr(Z=z)\log\Pr(Z=z)\), with
\(0\log0=0\), is finite entropy in nats; conditional entropy is
the average entropy of the conditional distributions.
\end{proposition}

\begin{proof}
An inverse recovers a unique antecedent, which requires injectivity;
conversely, injectivity makes it possible to assign to each expressed pair its unique
antecedent. Since \(q_k(Z)\) is a deterministic function of \(Z\),
grouping the sum defining \(\mathsf H(Z)\) by fibers of \(q_k\) yields the
chain rule. If the pair \((q_k(Z),M_k(Z))\) distinguishes the support,
each corresponding conditional distribution is concentrated on a single
state, and its entropy is zero.
\end{proof}

Not publishing a coordinate and erasing the record that permits its recovery are
distinct operations. This result does not identify entropy with the cardinality
of the continuum or introduce a thermal interpretation: its domain, reader, and
distribution are precisely those just specified.

\subsubsection{Equivariance of the measure under tree symmetries}
\label{act21:sec:ch-naturalidad-medida}

\begin{proposition}[Joint transport of the tree and its measure]
\label{act21:prop:ch-naturalidad-medida}
Let \(g=(g_k)_{k\geq0}\) be a collection of bijections \(g_k:H_k\to H_k\)
that commutes with truncations, fixes the root, and bijectively transports
admissible emissions while preserving their multiplicities. Then
\[
 d(g_k\zeta)=d(\zeta),\qquad
 (g_k)_*\mu_k=\mu_k,\qquad (g_\infty)_*\mu=\mu.
\]
More generally, if \(\mu_\zeta\) is the law of continuations of a
history \(\zeta\in H_k\), then
\[
 (g_\infty)_*\mu_\zeta=\mu_{g_k\zeta}.
\]
\end{proposition}

\begin{proof}
The bijection of emissions preserves their number with multiplicities. The
uniform seed distribution is likewise invariant. Each path
and its image therefore receive the same product of local weights
\(1/d(\zeta)\). Equality of measures follows at every level and on
every cylinder. Compatibility with truncations defines \(g_\infty\),
and uniqueness of the measure determined by cylinders proves equality
in the limit. The same argument, beginning at \(\zeta\), proves the
assertion about continuations.
\end{proof}

This establishes the naturality of the tree and its measure: a coherent change
of representation transports the weights and does not require choosing them anew.
Preservation of emissions and multiplicities is an essential hypothesis of
the proposition, not a datum added after the construction.


\subsection{Transport by sheets and exact accounting of outputs}
First consider the space with orthonormal basis formed by the histories at level \(k\). The transport preserving the label of each prolongation is

\[
\mathcal U_ke_\zeta=
\sum_{\varepsilon\ \mathrm{admissible\ at}\ \zeta}
\sqrt{p_{\rm can}(\varepsilon\mid\zeta)}\,e_{\zeta\varepsilon}.
\]
The successor fibers of distinct antecedents are disjoint. The normalization of \(p_{\rm can}\) therefore proves \(\mathcal U_k^*\mathcal U_k=I\). The representation by masses of histories is related to this one by the diagonal change of basis incorporating \(\sqrt{\mu(\zeta)}\).

The sheet label is part of the complete record of a history. We write
\[
\widehat\zeta=((\zeta_0,s_0),\ldots,(\zeta_k,s_k)),\qquad
s_j\in\{\Sigma,\Pi\},
\]
where \(s_j\) is the active sheet at step \(j\). The prolongation adds the emission and the sheet \(s_{k+1}\) determined by the operational TPK selector, retaining the prefix in full, including the initial sheet. In the basis of these lifted histories,
\[
\mathcal U_ke_{\widehat\zeta}
=\sum_{\varepsilon\ \mathrm{admissible\ at}\ \widehat\zeta}
\sqrt{p_{\rm can}(\varepsilon\mid\widehat\zeta)}\,e_{\widehat\zeta\varepsilon},
\qquad
P_ke_{\widehat\zeta}=\mathbf1_{s_k=\Sigma}e_{\widehat\zeta},\quad
Q_k=I-P_k.
\]
Thus, \(P_k\) reads the current sheet, while the distinction between histories is maintained by their complete prefixes. Two initial histories that reach the same active sheet retain distinct successor fibers. The proof of isometry by disjoint collections of prolongations applies verbatim to this basis.
The tritic orientation is another coordinate of the record: the value \(0\) preserves the threshold reading on each sheet and does not add a third sheet to this decomposition. We now separate sheet preservation and sheet exchange:

\[
A_k=P_{k+1}\mathcal U_kP_k+Q_{k+1}\mathcal U_kQ_k,
\qquad
\eta_k=Q_{k+1}\mathcal U_kP_k+P_{k+1}\mathcal U_kQ_k.
\]
\begin{proposition}[Gram identity by sheets]
\label{act21:prop:continuo-gram-hojas}
For a linear transport \(U\), whether or not it is isometric, the preceding decomposition satisfies

\[
A^*A+\eta^*\eta=PU^*UP+QU^*UQ.
\]
In particular, for \(U=\mathcal U_k\),

\[
A_k^*A_k+\eta_k^*\eta_k=I.
\]
\end{proposition}
\begin{proof}
In \(A^*A\) and \(\eta^*\eta\), the mixed terms containing \(P_{k+1}Q_{k+1}\) vanish. Summing the remaining terms produces \(P_kU^*(P_{k+1}+Q_{k+1})UP_k\), and its analogue with \(Q_k\). The first identity follows from complementarity of the projectors; the second also uses \(U^*U=I\).
\end{proof}

The first identity is the general formula. Replacing its right-hand side by \(U^*U\) requires the cross-blocks of that Gram matrix to vanish. Isometric preservation of history transport provides that condition here; it must not be attributed automatically to a subsequent compression.

Write \(F_{0,0}=I\), \(F_{0,n}=A_{n-1}\cdots A_0\), and \(T_N=F_{0,N}^*F_{0,N}\).

\begin{theorem}[Finite decomposition with terminal residue]
\label{act21:thm:continuo-descomposicion-terminal}
For each horizon \(N\),

\[
I=\sum_{n=0}^{N-1}
F_{0,n}^*\eta_n^*\eta_nF_{0,n}+T_N.
\]
The sequence \(T_N\) is positive and decreasing, and has a positive strong limit \(T_\infty\). The sum of outputs converges strongly to \(I-T_\infty\).
\end{theorem}
\begin{proof}
Proposition~\ref{act21:prop:continuo-gram-hojas} gives \(T_n-T_{n+1}=F_{0,n}^*\eta_n^*\eta_nF_{0,n}\). The telescoping sum proves the finite equality and monotonicity. For each vector, the quadratic forms decrease to a limit; polarization defines a positive bounded form and, by the representation theorem for bounded forms, an operator \(T_\infty\). Convergence is strong because, for a positive operator \(S\leq I\), \(\|Sv\|^2\leq\langle v,Sv\rangle\); this is applied to \(S=T_N-T_\infty\).
\end{proof}

The terminal residue is retained until its extinction in the domain under consideration is proved. This distinction prevents a phase return, by itself, from becoming an assumption that all histories are exhausted.

\subsection{Oriented balance and cancellation under refinement}
A history carries its accumulated orientation \(o(h)\) and the counts of its visits to the two sheets. The corresponding local balance is

\[
b(h)=o(h)\bigl(N_\Sigma(h)-N_\Pi(h)\bigr),
\qquad
b^+=\frac{|b|+b}{2},\quad b^-=\frac{|b|-b}{2}.
\]
When a reader expresses the mean over the prolongations, \(b_c=Eb_f\) is defined. The orientation remains in the argument of \(b_f\); it is not removed before averaging.

\begin{proposition}[Exact cancellation defect]
\label{act21:prop:continuo-defecto-cancelacion}
The quantity

\[
\kappa_{\mathrm{can}}=
\frac{E|b_f|-|Eb_f|}{2}
\]
is nonnegative and satisfies

\[
E(b_f^+)=(b_c)^++\kappa_{\mathrm{can}},
\qquad
E(b_f^-)=(b_c)^-+\kappa_{\mathrm{can}}.
\]
\end{proposition}
\begin{proof}
The weighted triangle inequality gives \(|Eb_f|\leq E|b_f|\). Writing the positive and negative parts as \((|b|\pm b)/2\) yields both equalities.
\end{proof}

The information omitted when only the balance is expressed is the common cancellation of the opposing contributions. This quantity is determined from the same prolongations used by the preceding transport. The factor \(1/9\) appears only when the fiber contains nine prolongations with equal weights; conditional expectation is the rule that remains valid for general branching and weightings.

\subsection{Ternary incidence of the sheets and rectangular complex}
The two sheets determine the module \(V=\mathbb F_3e_\Sigma\oplus\mathbb F_3e_\Pi\). Its four projective directions are the lines generated by \(e_\Sigma,e_\Pi,e_\Sigma+e_\Pi,e_\Sigma-e_\Pi\). They give rise to six unordered pairs. The eight nonzero vectors are their oriented representatives. Thus, the counts \(4,6,8\) belong to distinct operations on the same sheet module.

On the six pair positions, define

\[
(Ba)_{\{L,L'\}}=a_L+a_{L'},
\qquad
s(q)=\sum_{\{L,L'\}}q_{\{L,L'\}}.
\]
Each direction occurs in three pairs; hence \(sB=0\) in \(\mathbb F_3\). The reader \(W(q)=\mathbf1_{s(q)=0}-1/3\) is invariant under \(q\mapsto q+Ba\). On the uniform ambient space \(\mathbb F_3^6\), its moments are

\[
\mathbb EW=0,\qquad
\mathbb EW^2=\frac29,\qquad
\mathbb EW^3=\frac2{27}.
\]
Indeed, \(s\) is surjective and each fiber contains \(3^5\) elements. The reader equals \(2/3\) with probability \(1/3\) and \(-1/3\) with probability \(2/3\). These means describe the uniform ambient space; the distribution induced by a subset of histories must be calculated with its own multiplicities, already retained in \(H_k\).

To fix the ports without losing multiplicities, at depth \(k\) the
labeled readings are retained:
\(\lambda_\Sigma(\zeta)=(\zeta,\Sigma)\) and
\(\lambda_\Pi(\zeta)=(\zeta,\Pi)\). Their sets are
\[
\mathscr P_k^\Sigma=\{\lambda_\Sigma(\zeta):\zeta\in H_k\},\qquad
\mathscr P_k^\Pi=\{\lambda_\Pi(\zeta):\zeta\in H_k\}.
\]
They are finite and nonempty because \(H_k\) is. Values, quotients, and residues
are expressed on these ports; equality of values does not identify their
labels. The product \(\mathscr P_k^\Sigma\times\mathscr P_k^\Pi\) is the
domain of the rectangular reading, within which the support of the
admissible incidences is recorded.

With this level fixed, write \(I=\mathscr P_k^\Sigma\),
\(J=\mathscr P_k^\Pi\), and \(A_m=\mathbb Z/3^m\mathbb Z\). The
lexicographic APP order determines the minimal surviving branch and, within it, the
reference ports \(i_0,j_0\). No external choice is added: changing
coordinates produces a canonically isomorphic complex. The maps

\[
\kappa(u,v)_{ij}=u_i-v_j,
\qquad
(\delta w)_{ij}=w_{ij}-w_{ij_0}-w_{i_0j}+w_{i_0j_0}
\]
define the following sequence, valid over the finite ring \(A_m\):

\[
0\longrightarrow A_m\longrightarrow
A_m^I\oplus A_m^J\mathrel{\mathop{\longrightarrow}^{\kappa}}
A_m^{I\times J}\mathrel{\mathop{\longrightarrow}^{\delta}}
A_m^{(I\setminus\{i_0\})\times(J\setminus\{j_0\})}
\longrightarrow0.
\]
\begin{proposition}[Exactness of rectangular incidence]
\label{act21:prop:continuo-exactitud-rectangular}
The preceding sequence is exact, with its first arrow given by the constant diagonal.
\end{proposition}
\begin{proof}
The equality \(u_i-v_j=0\) for every pair forces all coordinates of \(u\) and \(v\) to be the same constant. Direct substitution gives \(\delta\kappa=0\). If \(\delta w=0\), take \(u_i=w_{ij_0}\) and \(v_j=w_{i_0j_0}-w_{i_0j}\). Then \(u_i-v_j=w_{ij}\). Finally, any matrix in the last term is prolonged with zero reference row and column; its image under \(\delta\) is the initial matrix. The proof uses only additions and subtractions, and therefore does not require division by a unit of the ring.
\end{proof}

Reduction from \(A_{m+1}\) to \(A_m\), while keeping
\(k,\mathscr P_k^\Sigma,\mathscr P_k^\Pi,i_0,j_0\) fixed, commutes with both
maps. The reconstruction formulas are the same at every modular
depth. This is the arithmetic compatibility proved in this subsection. The
prolongation of histories \(k\to k+1\) preserves the prefix labels; its
maps between sets of ports are data distinct from the change of
ring \(m\to m+1\).

Those maps are obtained from deletion of emissions:
\[
\rho_k^\Sigma(\zeta\varepsilon,\Sigma)=(\zeta,\Sigma),\qquad
\rho_k^\Pi(\zeta\varepsilon,\Pi)=(\zeta,\Pi).
\]
The minimal surviving branch of the APP order just constructed has compatible prolongations. Its two labels define compatible reference ports at every level. Coordinate transport is precomposition:
\[
(L_ku)_{i'}=u_{\rho_k^\Sigma(i')},\quad
(L_kv)_{j'}=v_{\rho_k^\Pi(j')},\quad
(L_kw)_{i'j'}=w_{\rho_k^\Sigma(i'),\rho_k^\Pi(j')}.
\]
For the last module in the sequence, the matrix is first prolonged by zero on the reference row and column, this same formula is applied, and the result is restricted to the new complement of the reference ports. This rule also defines \(L_k\) when a new port has the marked port as its antecedent.

\begin{proposition}[Naturality of the rectangular complex under prolongation]
\label{act21:iv:rectangular-natural}
The preceding maps satisfy
\(L_k\kappa_k=\kappa_{k+1}L_k\) and
\(L_k\delta_k=\delta_{k+1}L_k\).
They also commute with coefficient reduction
\(A_{m+1}\to A_m\) and compose in accordance with prefix deletion.
\end{proposition}
\begin{proof}
The first identity follows from
\(u_{\rho_k^\Sigma(i')}-v_{\rho_k^\Pi(j')}\).
In the second, each of the four terms defining \(\delta\)
is transported by the same precomposition. Compatibility of the
marked ports identifies their rows and columns; when an antecedent is
the marked port, the four terms cancel and agree with
prolongation by zero. All formulas are sums and differences with
integer coefficients, so they commute with modular reduction.
The composition of precompositions is precomposition with the
composite prefix, proving the final assertion.
\end{proof}

\subsection{Memory complex and compatibility of incidences}
The same accumulated memory coordinate \(c_y\in\mathbb Z\), at a residence \(y\), defines a free integral complex. Its reduction to \(A_m\) uses the same integer record before reduction; therefore, successive depths are compatible. Its generators are \(f_y\) in degree \(2\), \(e_y,b_y,r_y\) in degree \(1\), and \(g_y\) in degree \(0\):

\[
\partial f_y=3c_y e_y+b_y,\qquad
\partial e_y=g_y,\qquad
\partial b_y=-3c_y g_y,\qquad
\partial r_y=0.
\]
The identity \(\partial^2=0\) is verified on \(f_y\) by \(3c_yg_y-3c_yg_y=0\), and is immediate on the other generators.

If a path \(\alpha:y\to y'\) updates the memory, its transport in the complex sends the like-named generators to their new representatives, except for

\[
\Psi_\alpha(b_y)=b_{y'}+3(c_{y'}-c_y)e_{y'}.
\]
\begin{proposition}[Naturality of the memory complex]
\label{act21:prop:continuo-naturalidad-memoria}
The identities \(\partial\Psi_\alpha=\Psi_\alpha\partial\) and \(\Psi_{\beta\alpha}=\Psi_\beta\Psi_\alpha\) hold. The maps extend linearly over the chosen coefficient ring.
\end{proposition}
\begin{proof}
On \(f_y\), the image of its boundary is \(3c_ye_{y'}+b_{y'}+3(c_{y'}-c_y)e_{y'} =3c_{y'}e_{y'}+b_{y'}\). On \(b_y\), the boundary of its image is \(-3c_{y'}g_{y'}+3(c_{y'}-c_y)g_{y'}=-3c_yg_{y'}\). The other two cases are immediate. Under concatenation, the memory differences sum telescopically, proving the second identity.
\end{proof}

The relation between two representatives of the same residence can also be written as an explicit boundary. For an integer coordinate \(v\), reduced in the same ring when appropriate, let

\[
z^-=r_y,\qquad z^+=r_y+3c_yv e_y+v b_y,
\qquad h_*=-vf_y.
\]
Then \(\partial h_*=z^--z^+\). Under a transport that also updates \(v\), the correction \(K_\alpha=(v_{y'}-v_y)f_{y'}\) satisfies \(K_{\beta\alpha}=K_\beta+\Psi_\beta K_\alpha\). The symbols \(K_\alpha\) in this local homotopy do not denote the dodecaphase seal \(K\): they belong to another domain, and their function is specified by the preceding identity. Their function can also be checked in the equality

\[
z^+_{y'}-\Psi_\alpha z^+_y=\partial K_\alpha.
\]
Both sides equal \((v_{y'}-v_y)(3c_{y'}e_{y'}+b_{y'})\).

\begin{proposition}[Homology and memory of the local complex]
\label{act21:prop:continuo-homologia-memoria}
Over \(\mathbb Z\), or after reduction to \(A_m\), one has

\[
H_0(\Gamma)=H_2(\Gamma)=0,
\qquad H_1(\Gamma)\cong A_m[r_y]
\]
in the modular realization, and \(H_1(\Gamma)\cong\mathbb Z[r_y]\) in the integral realization.
\end{proposition}
\Needspace{9\baselineskip}
\begin{proof}
Define the degree-\(1\) homotopy by \(h(g_y)=e_y\), \(h(b_y)=f_y\), and zero on the other generators. Let \(p\) be the projection onto the summand generated by \(r_y\), and \(\iota\) its inclusion. Verification on each generator gives

\[
\partial h+h\partial=I-\iota p.
\]
In particular, on \(b_y\) the terms \(3c_ye_y\) and \(-3c_ye_y\) cancel. The complex deformation retracts onto the module generated by \(r_y\), concentrated in degree \(1\), proving the assertion.
\end{proof}

The transport \(\Psi_\alpha\) is a shear of determinant \(1\). Thus, this complex explicitly preserves how \(c_y\) changes at chain level, but its local homology does not distinguish the values of that coordinate. The complete memory remains in the labels and transports of the source state; it is not recovered from \(H_1\) alone. A memory-dependent numerical torsion requires the specified assembly and determinant reading, in addition to this local complex.

\subsection{Assembly with terminal paths}
For a finite horizon, the residences and their prolongations form a path category. Each residence \(\nu\) is assigned the free module \(W(\nu)=A_m[\operatorname{Path}(\nu,\mathrm{Term})]\), which retains each terminal continuation as a generator. An initial path acts on \(W\) by precomposition. The complex \(\Gamma(\nu)\) of the preceding subsection acts covariantly through \(\Psi\).

The assembly uses both transports in the same bigraded module:

\[
\mathcal B_{q,r}=
\bigoplus_{\nu_0\to\cdots\to\nu_q}
W(\nu_q)\otimes_{A_m}\Gamma_r(\nu_0).
\]
The normalized path complex is used: degenerate chains containing identities are omitted. The finite horizon then bounds the number of strict prolongations and the horizontal degree. The horizontal boundary deletes an intermediate residence by composing the adjacent arrows. At the first endpoint it transports \(\Gamma\); at the last it precomposes \(W\). Alternating signs give the path boundary \(d_{\mathrm{bar}}\). The total differential is

\[
D=d_{\mathrm{bar}}+(-1)^q\partial_\Gamma.
\]
The face identities cancel the terms of \(d_{\mathrm{bar}}^2\) pairwise. Naturality in Proposition~\ref{act21:prop:continuo-naturalidad-memoria} implies that \(d_{\mathrm{bar}}\) commutes with \(\partial_\Gamma\); the factor \((-1)^q\) cancels the cross terms in the square. Therefore, \(D^2=0\). Compatibility is proved on path generators and is not introduced by a label asserting membership in a complex.

Terminal evaluation also does not require choosing a particular history. For each continuation \(h\), the effective composition of its updates is applied. The result retains the pair formed by \(h\) and its reading, with its multiplicity. If \(k<\ell<N\), the continuations decompose as a disjoint union of prefixes up to \(\ell\) and their continuations up to \(N\). Proposition~\ref{act21:prop:continuo-compatibilidad-afin} proves that direct evaluation and evaluation of those two parts produce the same reading. This yields the terminal naturality square.

The path modules, incidences, memory transport, and their differentials arise from the same prolongations. Their assembly does not consist in adding to a state five coordinates that would be retained by definition: their relations are constructed through sums, boundaries, and effective compositions, and are verified on their generators.

\subsection{Determinants, limit, and arithmetic reading}
At each residence, the free local complex \(\Gamma\) admits its graded determinant line

\[
\operatorname{Det}\Gamma=
\bigotimes_r\left(\bigwedge^{\mathrm{rk}\Gamma_r}\Gamma_r\right)^{(-1)^r}.
\]
The total assembly of the preceding subsection, in turn, has the modules \(\operatorname{Tot}_d\mathcal B=\bigoplus_{q+r=d}\mathcal B_{q,r}\). At finite horizon and over the same ring, they are free of finite rank and only finitely many of them are nonzero. Its line is
\[
\operatorname{Det}(\operatorname{Tot}\mathcal B)=
\bigotimes_{q,r}
\left(\bigwedge^{\operatorname{rk}\mathcal B_{q,r}}
\mathcal B_{q,r}\right)^{(-1)^{q+r}}.
\]
This formula follows from the direct-sum decomposition in each degree. It distinguishes the line of the local complex from that of the assembly with all paths; their bases are ordered by means of the retained labels. Modular reduction commutes with these exterior powers because the modules and their bases arise from the same free integral modules at fixed horizon.
The line is defined by the modules and their degrees, both over \(\mathbb Z\) and after reduction to \(A_m\); a negative exponent denotes the dual module. A nonzero numerical torsion additionally requires specifying the complex over a field, the bases, and the relevant homology identifications. To calculate Smith factors, one uses the integer matrix produced before reduction, not an arbitrary lift of a modular matrix. An integer square block with nonzero determinant is invertible over \(\mathbb Q\); its reduction is invertible over \(A_m\) only when that determinant is a unit, that is, when it is not divisible by \(3\). Refinement of memories and incidences makes it possible to formulate and verify the stability of the relevant invariant factors; these conditions do not follow from the existence of the determinant line.

Passage to the limit preserves the compatible projective systems of residues, the histories, and their multiplicities. The squares proved in the preceding subsections remain commutative because their maps are given by the same formulas at each level. For an Archimedean coordinate, the additional step consists in publishing nonempty cylinder intervals, compatible under inclusion and with diameter tending to zero. Their nested compact closures have a one-point intersection: existence follows from compactness and uniqueness from the diameter bound. That point is the value of the limit reading. If the intervals are half-open, the point may belong only to their closures; the boundary representation rule then determines its positional notation. The value does not replace the history that produced it. The other representation preserves the incidences of the same state. The two targets share their origin without thereby receiving the same information.

The coinductive regions of Section~\ref{sec:generacion}, the dodecaphase
record of Section~\ref{sec:k-registro}, and the exceptional incidence
of Section~\ref{act21:chap:despliegue-selector-k} are expressed from that substrate with their
respective maps. The Hilbert-space realization subsequently retains the phase
labels and specifies its own unital inclusion; duality and the screens
act on the codomains constructed there. The joint operations assembled here
transmit histories, sheets, orientation, and memory to those
representations. Identifying a specific functional with an external analytic
form in turn requires writing its action on the generators and verifying
their inner products: the conservation identities of the present chapter
do not replace that composition.

\subsection{Mathematical scope of the joint construction}
The constructions of state, tree, and measure; the correlative operations; and
the terminal assembly with naturality and limit are developed in full
in this article. In this chapter their balance, incidence, and memory
identities have been formulated and proved on the same
generators. The five constructions thus remain operations of a
single tree. Each specialization keeps its exact hypotheses visible:
isometry of transport, canonical weighting law, treatment of the terminal
residue, and nonsingularity of the determinant block, where applicable.

The preceding equations prove the memory cocycles, weighted cancellation,
rectangular exactness, and compatibility of the memory complex at every
finite horizon, as well as its compatible passage to the limit.
A subsequent analytic identification must compose its operator on these
generators and verify the additional properties of its codomain; it does not
alter or replace the joint construction proved here.

\iffalse
% Preliminary formulation retained outside the assembly; superseded by the
% typed and verified version loaded at the end of this block.
\subsection{Joint emergence of the five constructions}
\label{act21:sec:cinco-construcciones-continuo-ch}

The preceding operations are now assembled without turning them into five branches.
For a surviving history \(h\in X_k\), a horizon \(N>k\), and a
modular depth \(m\geq1\), let the common base record be
\[
 \mathcal B_{k,N}(h):=
 \bigl(\mathscr T_{k,N}(h),
       (\varepsilon_\alpha,s_\alpha,o_\alpha,
        r_{\Sigma,\alpha},q_{\Sigma,\alpha},
        r_{\Pi,\alpha},q_{\Pi,\alpha},
        \operatorname{Carr}_\alpha,
        \kappa_\alpha,\partial_\alpha,
        \operatorname{Surv}_\alpha)_\alpha;
       \operatorname{src},\operatorname{tgt},\circ,\rho\bigr),
 \tag{C1}
 \label{act21:eq:base-comun-registros}
\]
where \(\mathscr T_{k,N}(h)\) is the subtree of prolongations of \(h\)
up to horizon \(N\). Each component of the record is an output of the
APP divisions, the TRIT orientation, or the TPK updates described
in the preceding sections; none arises from a classical problem chosen
afterward.

On this same base, the five structured objects are defined:
\[
\begin{aligned}
 \mathcal F_{\rm sp}(\mathcal B_{k,N})
   &:=(\ell^2(X_j,\mu_j),\mathcal U_j,A_j,\eta_j,T_j)_{k\leq j\leq N},\\
 \mathcal F_{\rm sol}(\mathcal B_{k,N})
   &:=(b,\kappa_{\rm can},C_\alpha,\kappa_\alpha)_{\alpha\subset\mathscr T_{k,N}},\\
 \mathcal F_{\rm gau}(\mathcal B_{k,N})
   &:=(I_j,J_j,B_j,\kappa_j,\delta_j)_{k\leq j\leq N},\\
 \mathcal F_{\rm coh}^{(m)}(\mathcal B_{k,N})
   &:=(\Gamma_y\otimes A_m,\partial,\Psi_\alpha,H_\bullet)_{y,\alpha\subset\mathscr T_{k,N}},\\
 \mathcal F_{\rm det}^{(m)}(\mathcal B_{k,N})
   &:=(\operatorname{Tot}\mathcal B,\operatorname{Det}(\operatorname{Tot}\mathcal B))_{k,N,m}.
\end{aligned}
 \tag{C2}
\]
The discriminants \({\rm sp},{\rm sol},{\rm gau},{\rm coh},{\rm det}\)
name, respectively, the spectral, solenoidal, gauge,
cohomological, and determinant readings of a single object. They are not input domains.
If \(p_T\) forgets the structure added by the reader \(T\) and returns the
record \(\mathcal B_{k,N}(h)\), we write
\(\widetilde{\mathcal F}_T=(\mathcal B_{k,N},\mathcal F_T,p_T)\). The
correlative output is the fiber product
\[
\begin{aligned}
 \mathcal J_{k,N}^{(m)}(h):={}&
 \widetilde{\mathcal F}_{\rm sp}
 \mathop{\times}_{\mathcal B_{k,N}(h)}
 \widetilde{\mathcal F}_{\rm sol}
 \mathop{\times}_{\mathcal B_{k,N}(h)}
 \widetilde{\mathcal F}_{\rm gau}\\
 &\mathop{\times}_{\mathcal B_{k,N}(h)}
 \widetilde{\mathcal F}_{\rm coh}^{(m)}
 \mathop{\times}_{\mathcal B_{k,N}(h)}
 \widetilde{\mathcal F}_{\rm det}^{(m)}.
\end{aligned}
 \tag{C3}
 \label{act21:eq:constructor-correlativo}
\]
Equality of the five projections onto \(\mathcal B_{k,N}(h)\) requires
preservation of the same tree, the same edges, and the same record; hence
\eqref{act21:eq:constructor-correlativo} is not a Cartesian product of five
objects assembled retrospectively.

Deletion of the last layer of the tree induces
\(\operatorname{Res}_{N+1,N}\). The preceding propositions give, on the
generators, the squares
\[
\begin{gathered}
 \operatorname{Res}_{N+1,N}\mathcal F_T(\mathcal B_{k,N+1})
 =\mathcal F_T(\mathcal B_{k,N}),
 \qquad T\in\{{\rm sp},{\rm sol},{\rm gau},{\rm coh}\},\\
 \operatorname{Det}(\operatorname{Tot}\mathcal B_{k,N+1})
 \simeq
 \operatorname{Det}(\operatorname{Tot}\mathcal K_{N+1})
 \otimes
 \operatorname{Det}(\operatorname{Tot}\mathcal B_{k,N}),\\
 (\partial\Psi_\alpha=\Psi_\alpha\partial),\qquad
 (L_k\kappa_k=\kappa_{k+1}L_k),\qquad
 (L_k\delta_k=\delta_{k+1}L_k).
\end{gathered}
 \tag{C4}
 \label{act21:eq:naturalidad-D}
\]
The change \(A_{m+1}\to A_m\) likewise commutes with \(\partial\), \(\Psi\),
\(\kappa\), \(\delta\), and the determinant lines. One therefore obtains
a single pro-system in \((N,m)\).

Let \(\mathcal X_\chi\) be the limit of compatible enriched histories. For
\(\xi\in\mathcal X_\chi\), denote by \(\mathfrak L_N(\xi)\) the ordered
prefix of its \(N\) complete emissions. The joint horizon fibers
are
\[
 \mathfrak J_{N,\bullet}^{(m)}
 :=\coprod_{h\in X_0}
 \{(\mathfrak L,J):\mathfrak L\text{ is a branch of length }N,
       \ J=\mathcal J_{0,N}^{(m)}(h)\}.
 \tag{C5}
 \label{act21:eq:fibras-conjuntas-horizonte}
\]
The mark \(\mathfrak L\) does not retain a tautological copy of \(\xi\): it is
formed by the emissions and updates produced at each edge. The
joint generator is
\[
\begin{aligned}
 \operatorname{Gen}_{\rm cont}:\mathcal X_\chi
   &\longrightarrow\varprojlim_{N,m}\mathfrak J_{N,\bullet}^{(m)},\\
 \xi&\longmapsto
   (\mathfrak L_N(\xi),\mathcal J_{0,N}^{(m)}(h_0))_{N,m},\\
 \mathcal C_{\rm cont}^{\rm disc}
   &:=\operatorname{im}(\operatorname{Gen}_{\rm cont}).
\end{aligned}
 \tag{C6}
 \label{act21:eq:continuo-discreto}
\]

\begin{theorem}[Correlative emergence of the five constructions]
\label{act21:thm:correlacion-cinco}
For every surviving history \(h\in X_k\), every horizon
\(N\geq k+1\), and every depth \(m\geq1\), the fiber product
\eqref{act21:eq:constructor-correlativo} exists and is nonempty. Its five
components emerge from the same APP--TRIT--TPK record, satisfy the
naturality squares \eqref{act21:eq:naturalidad-D}, the changes of base, and the
passage to the limit. The image \(\mathcal C_{\rm cont}^{\rm disc}\) contains five
inseparable aspects of the same continuum, not five continua or five
independent applications.
\end{theorem}

\begin{proof}
Prolongability of the tree provides at least one edge at each interior
vertex. The record of that edge determines the isometric transport and its
decomposition by sheets, the oriented cocycle, the incidence ports, the
memory complex, and its determinant line. In this way,
\(\mathcal F_{\rm sp}\), \(\mathcal F_{\rm sol}\),
\(\mathcal F_{\rm gau}\), \(\mathcal F_{\rm coh}^{(m)}\), and
\(\mathcal F_{\rm det}^{(m)}\) are constructed successively on the same base, proving the
nonemptiness of \eqref{act21:eq:constructor-correlativo}. None of the five
constructions appears as an input coordinate.

The affine compatibility proposition proves the composition of memory; the
Gram identity, rectangular exactness, and naturality of the complex
prove the first four lines of \eqref{act21:eq:naturalidad-D}. The
decomposition of the bicomplex by the last layer proves the multiplicative
identity of the determinant line. All these operations are calculated
on the same generators and commute with modular reduction. The fiber,
relation, number, orientation, carry, memory, boundary, and
route remain in \(\mathfrak L_N(\xi)\); the multisection, terminals,
and linear and probabilistic readers remain in
\(\mathcal J_{0,N}^{(m)}(h_0)\). By compatibility, the limit produces
\eqref{act21:eq:continuo-discreto} without disaggregation.
\end{proof}

The subsequent recognition of the single continuum constructed through its five
spectral, solenoidal, gauge, cohomological, and determinant realizations
belongs to subsequent conventional language. It neither selects its states nor
participates in its generation.
\fi

\subsection{Joint emergence of the five constructions}
\label{act21:sec:cinco-construcciones-continuo-ch}

For \(0\leq k\leq N\), let
\[
 \rho_{N,k}:=\rho_k\circ\rho_{k+1}\circ\cdots\circ\rho_{N-1},
 \qquad \rho_{k,k}:=\operatorname{id}_{H_k},
\]
and, for a surviving history \(\zeta\in H_k\), define
\[
 \operatorname{Desc}_{k,N}(\zeta)
 :=\{x\in H_N:\rho_{N,k}(x)=\zeta\}.
\]
The finite tree \(\mathscr T_{k,N}(\zeta)\) has as vertices
\[
 \coprod_{j=k}^{N}\{y\in H_j:\rho_{j,k}(y)=\zeta\}
\]
and as edges only the TPK prolongations
\(a:y\to z\) with \(z\in\rho_j^{-1}(y)\). Therefore, neither its vertices nor
its prolongations arise from an auxiliary tree.

For \(x\in\operatorname{Desc}_{k,N}(\zeta)\), the marked branch
\(\lambda_{\zeta,x}\) is the sequence
\[
 \zeta=x_k\longrightarrow x_{k+1}\longrightarrow\cdots
 \longrightarrow x_N=x
\]
determined by the truncations of \(x\). The common base record is
\[
\begin{aligned}
 \mathscr R_{k,N}(\zeta;x):=
 \bigl(&\mathscr T_{k,N}(\zeta),\lambda_{\zeta,x};\\
 &\bigl(\varepsilon_a,s_a,o_a,
   (\delta_{\diamond,a},r_{\diamond,a},q_{\diamond,a})
       _{\diamond\in\{\Sigma,\Pi\}},
   \operatorname{Carr}_a,C_a,t_a,\operatorname{Fr}_a,
   \operatorname{Route}_a,\operatorname{Mem}_a,
   \operatorname{Surv}_a\bigr)_a;\\
 &\operatorname{src},\operatorname{tgt},\circ,
   (\rho_j)_{k\leq j<N}\bigr),
\end{aligned}
\tag{C1}
\label{act21:eq:base-comun-registros}
\]
where
\[
 U_a(m_0)=C_am_0+t_a,
 \qquad
 \delta_{\diamond,a}=r_{\diamond,a}+9q_{\diamond,a}.
\]
The symbol \(m_0\) in this formula is a memory coordinate; it is not the
modular depth \(m\) used below. The true base space is
\[
 \mathfrak R_{k,N}(\zeta)
 :=\{\mathscr R_{k,N}(\zeta;x):
       x\in\operatorname{Desc}_{k,N}(\zeta)\}.
\]
The marked branch is not an added input: it is the sequence of edges
actually produced that ends at \(x\). The remainder of the record simultaneously
preserves the lateral fibers of the tree; therefore, the mark does not replace
the branching with an isolated trajectory.

\Needspace{15\baselineskip}
Set \(R_m:=\mathbb Z/3^m\mathbb Z\), and let
\(V_j(\zeta):=\operatorname{Desc}_{k,j}(\zeta)\). The transport, the
sheet projectors, and the terminal residue are restricted to the subtree by
\[
 \mathcal U_j^\zeta:=
 \mathcal U_j\!\upharpoonright_{\ell^2(V_j(\zeta),\mu_j)},
 \qquad P_j^\zeta:=P_j\!\upharpoonright_{V_j(\zeta)},
 \qquad Q_j^\zeta:=I-P_j^\zeta,
\]
\[
 A_j^\zeta:=P_{j+1}^\zeta\mathcal U_j^\zeta P_j^\zeta
             +Q_{j+1}^\zeta\mathcal U_j^\zeta Q_j^\zeta,
 \qquad
 \eta_j^\zeta:=Q_{j+1}^\zeta\mathcal U_j^\zeta P_j^\zeta
             +P_{j+1}^\zeta\mathcal U_j^\zeta Q_j^\zeta.
\]
For \(k\leq n\leq N\), let
\[
 F_{k,k}^\zeta:=I,
 \qquad
 F_{k,n}^\zeta:=A_{n-1}^\zeta\cdots A_k^\zeta,
 \qquad
 T_{k,n}^\zeta:=(F_{k,n}^\zeta)^*F_{k,n}^\zeta.
\]
The proper ports of the same subtree are
\[
 \mathscr P_j^\Sigma(\zeta):=
 \{(y,\Sigma):y\in V_j(\zeta)\},
 \qquad
 \mathscr P_j^\Pi(\zeta):=
 \{(y,\Pi):y\in V_j(\zeta)\};
\]
we denote by \(B_j^\zeta,\kappa_j^\zeta,\delta_j^\zeta\) the restrictions
of the incidence and of the rectangular complex to those ports, and by
\(L_j^\zeta\) the precomposition induced by prefix deletion.

To define the path reader without abbreviations, let
\[
 W_\ell^\zeta(y):=
 R_m[\operatorname{Path}_{\mathscr T_{k,\ell}(\zeta)}
          (y,\operatorname{Term}_\ell)]
\]
and set
\[
 \mathscr D_{q,r;k,\ell}^{(m)}[\mathscr R]
 :=\bigoplus_{y_0\to\cdots\to y_q\ \text{in}\
                  \mathscr T_{k,\ell}(\zeta)}
 W_\ell^\zeta(y_q)\otimes_{R_m}
 \bigl(\Gamma_r(y_0)\otimes_{\mathbb Z}R_m\bigr),
 \qquad
 \operatorname{Tot}_d\mathscr D
 :=\bigoplus_{q+r=d}\mathscr D_{q,r;k,\ell}^{(m)}[\mathscr R].
\]
The formulas already proved produce, on each
\(\mathscr R\in\mathfrak R_{k,N}(\zeta)\), the five objects
\[
\begin{aligned}
 \mathcal F_{\rm sp}[\mathscr R]
 &:=(\ell^2(V_j(\zeta),\mu_j),
       \mathcal U_j^\zeta,A_j^\zeta,\eta_j^\zeta,
       T_{k,j}^\zeta)_{k\leq j\leq N},\\
 \mathcal F_{\rm sol}[\mathscr R]
 &:=(b,\kappa_{\rm can},C_\alpha,\kappa_\alpha)
      _{\alpha\subset\mathscr T_{k,N}(\zeta)},\\
 \mathcal F_{\rm gau}^{(m)}[\mathscr R]
 &:=(\mathscr P_j^\Sigma(\zeta),\mathscr P_j^\Pi(\zeta),
       B_j^\zeta,\kappa_j^\zeta,\delta_j^\zeta;R_m)_{k\leq j\leq N},\\
 \mathcal F_{\rm coh}^{(m)}[\mathscr R]
 &:=(\Gamma_y\otimes R_m,\partial,\Psi_\alpha,H_\bullet)
      _{y,\alpha\subset\mathscr T_{k,N}(\zeta)},\\
 \mathcal F_{\rm det}^{(m)}[\mathscr R]
 &:=(\operatorname{Det}(\operatorname{Tot}
       \mathscr D_{\bullet,\bullet;k,\ell}^{(m)}[\mathscr R]))
       _{k\leq\ell\leq N}.
\end{aligned}
\tag{C2}
\label{act21:eq:cinco-lectores-comunes}
\]
Here \(\mathscr D_{\bullet,\bullet;k,\ell}^{(m)}[\mathscr R]\) is the normalized
bicomplex
of paths and memory from the preceding sections, restricted to the subtree of
horizon \(\ell\) and reduced over \(R_m\). The last reader preserves the
ordered collection of determinant lines up to horizon \(N\); in this
way, its truncation is a canonical projection and does not presuppose a
nonexistent determinant factorization. The law \(\mu\) is the canonical law
of the census, and the ports are fixed by the lexicographically minimal branch
of the APP order. Therefore, none of the five readers receives an external weighting
or an external port.

For
\(T\in\mathfrak T:=\{{\rm sp},{\rm sol},{\rm gau},{\rm coh},{\rm det}\}\),
define the typed total space
\[
 \widetilde{\mathcal F}_{T;k,N}^{(m)}(\zeta)
 :=\coprod_{\mathscr R\in\mathfrak R_{k,N}(\zeta)}
   \{\mathcal F_T^{(m)}[\mathscr R]\}_{\mathscr R},
 \qquad
 p_T:\widetilde{\mathcal F}_{T;k,N}^{(m)}(\zeta)
      \longrightarrow\mathfrak R_{k,N}(\zeta),
\]
where
\(\mathcal F_{\rm sp}^{(m)}:=\mathcal F_{\rm sp}\) and
\(\mathcal F_{\rm sol}^{(m)}:=\mathcal F_{\rm sol}\); that is, the
superscript \(m\) is inert for those two readers.
The projection \(p_T\) forgets only the reader structure and preserves
its source record. The correct fiber product is
\[
\begin{aligned}
 \mathcal J_{k,N}^{(m)}(\zeta):={}&
 \widetilde{\mathcal F}_{\rm sp;k,N}^{(m)}(\zeta)
 \mathop{\times}_{\mathfrak R_{k,N}(\zeta)}
 \widetilde{\mathcal F}_{\rm sol;k,N}^{(m)}(\zeta)
 \mathop{\times}_{\mathfrak R_{k,N}(\zeta)}
 \widetilde{\mathcal F}_{\rm gau;k,N}^{(m)}(\zeta)\\
 &\mathop{\times}_{\mathfrak R_{k,N}(\zeta)}
 \widetilde{\mathcal F}_{\rm coh;k,N}^{(m)}(\zeta)
 \mathop{\times}_{\mathfrak R_{k,N}(\zeta)}
 \widetilde{\mathcal F}_{\rm det;k,N}^{(m)}(\zeta).
\end{aligned}
\tag{C3}
\label{act21:eq:constructor-correlativo}
\]
The equality of the five projections forces preservation of the same tree, the
same edges, the same produced branch, and the same record. For each
\(\mathscr R\), the five constructions determine the element
\[
 s_{k,N}^{(m)}(\mathscr R)
 :=(\mathcal F_T^{(m)}[\mathscr R])_{T\in\mathfrak T}
 \in\mathcal J_{k,N}^{(m)}(\zeta).
\]
The existence of this element arises from the operations of
\eqref{act21:eq:cinco-lectores-comunes}, not from the
fiber-product notation.

Deletion of the last layer induces
\[
 \operatorname{res}^{\mathfrak R}_{N+1,N}
 \bigl(\mathscr R_{k,N+1}(\zeta;x)\bigr)
 :=\mathscr R_{k,N}(\zeta;\rho_N(x)).
\]
The restriction of each operator to the retained generators defines
\[
 r^m_{N+1,N}:\mathcal J_{k,N+1}^{(m)}(\zeta)
 \longrightarrow\mathcal J_{k,N}^{(m)}(\zeta).
\]
In the determinantal reader, \(r^m_{N+1,N}\) deletes only the last
coordinate of the collection of lines in
\eqref{act21:eq:cinco-lectores-comunes}. The reduction
\(R_{m+1}\twoheadrightarrow R_m\) defines, by base change,
\[
 q^N_{m+1,m}:\mathcal J_{k,N}^{(m+1)}(\zeta)
 \longrightarrow\mathcal J_{k,N}^{(m)}(\zeta).
\]
On each generator, the following identities hold:
\[
\begin{gathered}
 r^m_{N+1,N}q^{N+1}_{m+1,m}
 =q^N_{m+1,m}r^{m+1}_{N+1,N},\\
 \partial\Psi_\alpha=\Psi_\alpha\partial,
 \qquad L_j^\zeta\kappa_j^\zeta
       =\kappa_{j+1}^\zeta L_j^\zeta,
 \qquad L_j^\zeta\delta_j^\zeta
       =\delta_{j+1}^\zeta L_j^\zeta,
\end{gathered}
\tag{C4}
\label{act21:eq:naturalidad-D}
\]
and the compositions of the maps \(r\) and \(q\) satisfy the functorial
identities. The first equality holds because deleting the last layer and
reducing the coefficients act on distinct coordinates. The other three
are the naturality identities proved above. A single
pro-system in \((N,m)\) is obtained, not five juxtaposed limits.

At the root, define
\[
 \mathfrak J_{N,\bullet}^{(m)}:=\mathcal J_{0,N}^{(m)}(\ast).
\tag{C5}
\label{act21:eq:fibras-conjuntas-horizonte}
\]
The joint discrete structure of the continuum is the limit of those
correlative objects:
\[
 \boxed{
 \mathcal C_{\rm cont}^{\rm disc}
 :=\varprojlim_{\substack{N\geq1\\m\geq1}}
   \bigl(\mathfrak J_{N,\bullet}^{(m)},
         r^m_{N+1,N},q^N_{m+1,m}\bigr).}
\tag{C6}
\label{act21:eq:continuo-discreto}
\]
This is a definition by the pro-system generated by the full tree; it is not
the image or the graph of a map that retains the input state as its
first coordinate.

The records possess a produced terminal mark. Their compatible projections
induce
\[
 \Pi:\mathcal C_{\rm cont}^{\rm disc}\longrightarrow H_\infty^{\rm enr}.
\]
Conversely, a complete history determines, through its already emitted prefixes,
the compatible sections \(s_{0,N}^{(m)}\). Thus, after having
defined the continuum, one obtains
\[
 \operatorname{Gen}_{\rm cont}:H_\infty^{\rm enr}
 \longrightarrow\mathcal C_{\rm cont}^{\rm disc},
 \qquad
 \Pi\circ\operatorname{Gen}_{\rm cont}
 =\operatorname{id}_{H_\infty^{\rm enr}}.
\tag{C7}
\label{act21:eq:seccion-historias-continuo}
\]
This section proves recoverability; it is not used to define
\(\mathcal C_{\rm cont}^{\rm disc}\).

\begin{theorem}[Correlative emergence of the five constructions]
\label{act21:thm:correlacion-cinco}
For every surviving history \(\zeta\in H_k\), every \(N>k\), every
\(m\geq1\), and every \(\mathscr R\in\mathfrak R_{k,N}(\zeta)\), the five
constructions of \eqref{act21:eq:cinco-lectores-comunes} are defined on the
same record, and \eqref{act21:eq:constructor-correlativo} contains
\(s_{k,N}^{(m)}(\mathscr R)\). Their truncations, modular reductions, and
transports satisfy \eqref{act21:eq:naturalidad-D}. Furthermore,
\(\mathcal C_{\rm cont}^{\rm disc}\neq\varnothing\), the map \(\Pi\)
is surjective and is split by \(\operatorname{Gen}_{\rm cont}\). The
five components are inseparable aspects of the same continuum, not five
domains or five independent maps.
\end{theorem}

\begin{proof}
The survival that defines \(H_k\) and the TPK prolongation law give
\(\rho_k^{-1}(\zeta)\neq\varnothing\) for each vertex. Because the fibers are
finite, Kőnig's lemma yields \(H_\infty^{\rm enr}\neq\varnothing\).

For a fixed record \(\mathscr R\), the transport of histories and sheets produces
\(\mathcal F_{\rm sp}[\mathscr R]\); affine composition and the exact
cancellation defect produce \(\mathcal F_{\rm sol}[\mathscr R]\); the incidences
and the rectangular sequence produce \(\mathcal F_{\rm gau}^{(m)}[\mathscr R]\);
the memory complex and its transport produce
\(\mathcal F_{\rm coh}^{(m)}[\mathscr R]\); and the lines of the normalized
path assembly produce \(\mathcal F_{\rm det}^{(m)}[\mathscr R]\).
Each construction acts on the same generated edges, not on an
assumed subdomain or on the abstract ambient space of \(729\) words.

The affine, transport, incidence, and memory compatibilities prove
\eqref{act21:eq:naturalidad-D}. Truncating the determinant collection of
\eqref{act21:eq:cinco-lectores-comunes} commutes by
definition with \(r\); its base change commutes with \(q\). The sections
\(s_{0,N}^{(m)}\) are compatible and thus produce
\(\operatorname{Gen}_{\rm cont}\). The identity
\eqref{act21:eq:seccion-historias-continuo} proves simultaneously the
nonemptiness of the limit, the surjectivity of \(\Pi\), and the recoverability of each
history without introducing it as a defining coordinate.

The fiber, relation, number, orientation, carry, memory,
boundary, and route remain in the common record; the multisection, the
terminals, and the linear and probabilistic readers remain in the five
correlative constructions. Therefore, passage to the limit preserves the closure of
dependencies and does not disaggregate the continuum.
\end{proof}

The subsequent recognition of the single continuum constructed through
spectral, solenoidal, gauge, cohomological, or determinant language belongs to the
subsequent conventional comparison. It neither selects its states nor intervenes in
its APP--TRIT--TPK generation.

